% GENERATED FILE. Edit canonical lessons, then run npm run generate:books.
% canonical-source-hash: fnv1a64:27974f1340ed82f9
% canonical-lessons: RU-C267-ya-byl-rad, RU-C267-mechtat, RU-C267-review

\chapter{Glad, Tired, Hopeful}
\label{ch:ru-glad-tired-hopeful}
\hlchaptermodality{267}

\begin{chapteropening}
\textbf{By the end of this chapter:} \emph{I can say how I felt with \ru{я} \ru{был} \ru{рад}, \ru{я} \ru{устал} and \ru{мне} \ru{было} \ru{скучно}, and what I hope for and dream of with \ru{я} \ru{надеюсь}, \ru{что} and \ru{я} \ru{мечтаю}.}

Complete RU-C267-review: how a trip felt, in the past, and one hope and one dream for the next one.
\end{chapteropening}

\section[\texorpdfstring{\ru{я} \ru{был} \ru{рад}, \ru{я} … (ya byl rad, ya …)}{я был рад, я … (ya byl rad, ya …)}]{\ru{я} \ru{был} \ru{рад}, \ru{я} \ru{устал}, \ru{мне} \ru{было} \ru{скучно} — how it felt}

\label{lesson:RU-C267-ya-byl-rad}



\begin{quote}
Say \textquotedblleft{}joy\textquotedblright{}, \textquotedblleft{}happy\textquotedblright{} and \textquotedblleft{}sad\textquotedblright{}. (\emph{Rádost. Schastlívyy. Grústnyy.}) Then \textquotedblleft{}I've missed the train\textquotedblright{}. (\emph{Ya opozdál na póyezd.}) A trip has problems, and it has feelings.
\end{quote}

\begin{grammarlens}[title={Grammar Lens: \ru{рад}, glad}]
\textbf{\ru{рад}} — \emph{rad} — \textquotedblleft{}glad\textquotedblright{}. It is a short adjective: it does not stand before a noun, only after the person, and it agrees: \textbf{\ru{рад}} for a man, \textbf{\ru{рада}} (\emph{ráda}) for a woman, \textbf{\ru{рады}} (\emph{rády}) for more than one.

\begin{quote}
\textbf{\ru{Я} \ru{рад}, \ru{что} \ru{ты} \ru{позвонил}.} — \emph{ya rad, shto ty pozvoníl} — I'm glad you rang. \textbf{\ru{Я} \ru{была} \ru{рада}, \ru{что} \ru{ты} \ru{позвонил}.} — \emph{ya bylá ráda, shto ty pozvoníl} — I was glad you rang.
\end{quote}

In the present there is no verb; in the past, \textbf{\ru{был}} agrees with it.
\end{grammarlens}

\begin{grammarlens}[title={Grammar Lens: \ru{я} \ru{устал}, I'm tired}]
\begin{quote}
\textbf{\ru{Я} \ru{устал}.} — \emph{ya ustál} — I'm tired. A woman: \textbf{\ru{Я} \ru{устала}.}
\end{quote}

This is the perfective past of \emph{\ru{устать}}, the partner of \emph{\ru{уставать}}: \textquotedblleft{}I have got tired\textquotedblright{}. The result is the state you are in now, so a past form says \textquotedblleft{}I am tired\textquotedblright{} today. \emph{\ru{Усталый}} is the adjective for a tired face or a tired voice.
\end{grammarlens}

\begin{grammarlens}[title={Grammar Lens: \ru{мне} \ru{было} \ru{скучно}}]
How a thing felt goes on the \textquotedblleft{}to me\textquotedblright{} form, with an \emph{-\ru{о}} word and \textbf{\ru{было}}:

\begin{quote}
\textbf{\ru{На} \ru{концерте} \ru{мне} \ru{было} \ru{скучно}.} — \emph{na kontsérte mne býlo skúchno} — I was bored at the concert. \textbf{\ru{Мне} \ru{было} \ru{интересно} \ru{и} \ru{весело}.} — \emph{mne býlo interésno i véselo} — It was interesting, and I had fun.
\end{quote}

\emph{\ru{Скучно}} is said \emph{skúshno}: the \emph{\ru{ч}} is said as \emph{\ru{ш}} in this word.
\end{grammarlens}

\begin{cousinweb}
\textbf{\ru{рад}} goes back to Proto-Slavic \emph{*radŭ}, \textquotedblleft{}glad\textquotedblright{} (EtymDB 2.1). \textbf{\ru{Радость}}, \textquotedblleft{}joy\textquotedblright{}, is the same root with a noun ending: the gladness itself.
\end{cousinweb}

\subsection*{Your turn}
\begin{itemize}
\raggedright
  \item \emph{Say it:} \emph{Ya rad, shto ty pozvoníl}, then a woman, in the past: \emph{Ya bylá ráda, shto ty pozvoníl.}
  \item \emph{Say it:} \emph{Ya ustál}, then \emph{Ya ustála.}
  \item \emph{Say it:} \emph{Na kontsérte mne býlo skúchno.}
  \item \emph{Answer:} \emph{Tebé býlo interésno?} Answer with \emph{mne býlo} and one \emph{-o} word
\end{itemize}

\begin{tcolorbox}[breakable,colback=teal!4,colframe=teal!35!black,title={Before you move on}]
\textquotedblleft{}She was glad\textquotedblright{}? (\emph{\textbf{\ru{Она} \ru{была} \ru{рада}.}}) \textquotedblleft{}I'm tired\textquotedblright{}, in the past form? (\emph{\textbf{\ru{Я} \ru{устал}}}, or \emph{\textbf{\ru{я} \ru{устала}}}.) \textquotedblleft{}I was bored\textquotedblright{}? (\emph{\textbf{\ru{Мне} \ru{было} \ru{скучно}.}})
\end{tcolorbox}
\section[\texorpdfstring{\ru{мечтать} (mechtát)}{мечтать (mechtát)}]{\ru{мечтать} (mechtát) — to dream of, to long to}

\label{lesson:RU-C267-mechtat}



\begin{quote}
Say \textquotedblleft{}I was glad\textquotedblright{} and \textquotedblleft{}I had fun\textquotedblright{}. (\emph{Ya byl rad. Mne býlo véselo.}) Then \textquotedblleft{}to hope\textquotedblright{}. (\emph{Nadéyatsya.}) A trip told in the past leaves something for the future: a hope, a dream.
\end{quote}

\subsection*{You'll want to know: \ru{мечтать}}
\textbf{\ru{мечтать}} — \emph{mechtát} — \textquotedblleft{}to dream of, to long to\textquotedblright{}. It takes an infinitive, as \emph{\ru{хотеть}} does:

\begin{quote}
\textbf{\ru{Я} \ru{мечтаю} \ru{путешествовать}.} — \emph{ya mechtáyu puteshéstvovat} — I dream of travelling.
\end{quote}

Or \textbf{\ru{о}} and a thing, with the thing's ending changed:

\begin{quote}
\textbf{\ru{Я} \ru{мечтаю} \ru{о} \ru{море}.} — \emph{ya mechtáyu o móre} — I dream of the sea.
\end{quote}

This is a waking dream, a wish. A dream at night is \textbf{\ru{сон}}, the word you have for sleep.

\begin{grammarlens}[title={Grammar Lens: \ru{я} \ru{надеюсь}, \ru{что} …}]
A dream may never come; a hope is about what may well happen. \textbf{\ru{Я} \ru{надеюсь}, \ru{что} …} takes a whole clause, and a hope about tomorrow puts it in the future:

\begin{quote}
\textbf{\ru{Я} \ru{надеюсь}, \ru{что} \ru{погода} \ru{будет} \ru{хорошая}.} — \emph{ya nadéyus, shto pogóda búdet khoróshaya} — I hope the weather will be good.
\end{quote}

\textbf{\ru{Будет}}, \textquotedblleft{}will be\textquotedblright{}, is the future of \emph{\ru{быть}}: the \emph{\ru{буду}} inside \emph{\ru{будущий}}. Short and spoken: \textbf{\ru{Надеюсь}!} — \textquotedblleft{}I hope so!\textquotedblright{}
\end{grammarlens}

\begin{cousinweb}
\textbf{\ru{мечтать}} is inherited from Proto-Slavic \emph{*mĭčĭtati}, and the noun \textbf{\ru{мечта}}, \textquotedblleft{}a dream, a wish\textquotedblright{}, from \emph{*mĭčĭta}, \textquotedblleft{}dream\textquotedblright{} (EtymDB 2.1). EtymDB derives that noun from a Proto-Slavic verb \emph{*mikati} and gives it no meaning; by the usual account it is the verb of blinking and flickering, so a dream was first a flicker before the eyes.
\end{cousinweb}

\subsection*{Your turn}
\begin{itemize}
\raggedright
  \item \emph{Say it:} \emph{Ya mechtáyu puteshéstvovat.}
  \item \emph{Say it:} \emph{Ya mechtáyu o móre.}
  \item \emph{Say it:} \emph{Ya nadéyus, shto pogóda búdet khoróshaya.}
  \item \emph{Answer:} \emph{O chyom ty mechtáyesh?} \textquotedblleft{}What do you dream of?\textquotedblright{} Answer with an infinitive or with \emph{o}
\end{itemize}

\begin{tcolorbox}[breakable,colback=teal!4,colframe=teal!35!black,title={Before you move on}]
\textquotedblleft{}I dream of travelling\textquotedblright{}? (\emph{\textbf{\ru{Я} \ru{мечтаю} \ru{путешествовать}.}}) A dream at night? (\emph{\textbf{\ru{Сон}.}}) \textquotedblleft{}I hope the weather will be good\textquotedblright{}? (\emph{\textbf{\ru{Я} \ru{надеюсь}, \ru{что} \ru{погода} \ru{будет} \ru{хорошая}.}})
\end{tcolorbox}
\section[(review)]{Review — how it felt, and what comes next}

\label{lesson:RU-C267-review}



\begin{quote}
Say \textquotedblleft{}I'm glad\textquotedblright{}, as a man and as a woman. (\emph{Ya rad. Ya ráda.}) Then \textquotedblleft{}I'm tired\textquotedblright{}. (\emph{Ya ustál. Ya ustála.})
\end{quote}

\begin{grammarlens}[title={What you've built}]
Two lessons, and an experience now has a feeling and a future. \textbf{\ru{Рад}}, \textbf{\ru{рада}}, \textbf{\ru{рады}}, with \emph{\ru{был}} in the past. \textbf{\ru{Я} \ru{устал}}, a past form for a state now. \textbf{\ru{Мне} \ru{было}} with an \emph{-\ru{о}} word for how a thing felt. And two ways forward: \textbf{\ru{я} \ru{надеюсь}, \ru{что}} … \emph{\ru{будет}}, and \textbf{\ru{я} \ru{мечтаю}} with an infinitive or \textbf{\ru{о}}.

A day at a museum, in three lines:

\begin{quote}
\textbf{\ru{В} \ru{музее} \ru{мне} \ru{было} \ru{интересно}. \ru{Вечером} \ru{я} \ru{устала}, \ru{но} \ru{я} \ru{была} \ru{рада}. \ru{Я} \ru{мечтаю} \ru{вернуться} \ru{туда}.}
\end{quote}

\emph{\ru{Вернуться}}, the partner of \emph{\ru{возвращаться}}, came with Anna's passport: \textquotedblleft{}I dream of going back there\textquotedblright{}.
\end{grammarlens}

\subsection*{Your turn}
\begin{itemize}
\raggedright
  \item \emph{Say it:} \emph{V muzéye mne býlo interésno.}
  \item \emph{Say it:} \emph{Ya opozdál na póyezd, no mne né bylo grústno.}
  \item \emph{Say it:} a late train as good news: \emph{Na vokzále býlo obyavléniye, i ya byl rad: póyezd opázdyval.}
  \item \emph{Say it:} \emph{Ya nadéyus, shto pogóda búdet khoróshaya.}
  \item \emph{Answer:} \emph{Kak býlo na kontsérte?} \textquotedblleft{}How was the concert?\textquotedblright{} Answer with \emph{mne býlo} and one feeling
\end{itemize}

\begin{tcolorbox}[breakable,colback=teal!4,colframe=teal!35!black,title={Before you move on}]
\textquotedblleft{}It was interesting for me\textquotedblright{}? (\emph{\textbf{\ru{Мне} \ru{было} \ru{интересно}.}}) \textquotedblleft{}I dream of the sea\textquotedblright{}? (\emph{\textbf{\ru{Я} \ru{мечтаю} \ru{о} \ru{море}.}}) A hope about tomorrow: which verb form? (\emph{\textbf{\ru{Будет}}}, the future.) And the silent sign after \emph{ob-} in \emph{obyavléniye}? (The \textbf{hard sign}, \emph{tvyórdyy znak}.)
\end{tcolorbox}
